\section{Відповіді на контрольні запитання}

\begin{description}
    \item[1. Що таке колекція в мові Java?] ~\\
    Колекція в Java — це спеціальний об'єкт, який служить контейнером і призначений для компактного зберігання, динамічного опрацювання та управління групою інших об'єктів. Колекції входять до складу \textit{Java Collections Framework} (пакет \texttt{java.util}).

    \item[2. Для чого використовуються колекції?] ~\\
    Колекції використовуються для ефективної роботи з динамічними наборами даних у пам'яті комп'ютера, коли заздалегідь невідома точна кількість елементів. Вони надають готові високоефективні інструменти (алгоритми) для швидкого додавання, видалення, сортування, фільтрації та пошуку об'єктів.

    \item[3. Чим колекція відрізняється від масиву?] ~\\
    Масив має строго фіксований розмір, який не можна змінити після створення, і може зберігати як примітиви, так і об'єкти. Колекція може динамічно розширюватися або стискатися в процесі роботи програми, має великий вбудований набір методів обробки і здатна зберігати \textbf{виключно об'єктні типи даних}.

    \item[4. Яке призначення інтерфейсу \texttt{List}?] ~\\
    Інтерфейс \texttt{List} призначений для представлення \textbf{впорядкованої послідовності елементів} (динамічного масиву або списку). Головні особливості списків: кожен елемент має свій чіткий індекс починаючи з нуля, а сама структура повністю допускає наявність однакових значень (дублікатів).

    \item[5. Які основні властивості \texttt{ArrayList}?] ~\\
    \texttt{ArrayList} — це реалізація інтерфейсу \texttt{List} на базі звичайного масиву, який автоматично збільшує свою ємність при заповненні. Він забезпечує дуже швидкий доступ до елементів за їхнім індексом за константний час $O(1)$, але працює повільніше при видаленні чи вставці елементів всередину списку через необхідність зсовування сусідніх комірок.

    \item[6. Для чого використовується метод \texttt{add()}?] ~\\
    Метод \texttt{add()} використовується для динамічного додавання нового елемента в кінець поточної колекції. У випадку інтерфейсу \texttt{List} цей метод завжди повертає \texttt{true} і збільшує розмір списку на одиницю.

    \item[7. Для чого використовується метод \texttt{get()}?] ~\\
    Метод \texttt{get(int index)} використовується виключно у списках \texttt{List} для миттєвого отримання посилання на об'єкт за його порядковим номером (індексом). Якщо вказати індекс, який виходить за межі фактичного розміру колекції, програма видасть помилку \texttt{IndexOutOfBoundsException}.

    \item[8. Яке призначення методу \texttt{set()}?] ~\\
    Метод \texttt{set(int index, E element)} призначений для заміни (перезапису) існуючого елемента списку за вказаним індексом на новий об'єкт. Метод повертає те старе значення, яке було замінене в цій комірці.

    \item[9. Чим відрізняються \texttt{remove()} та \texttt{clear()}?] ~\\
    Метод \texttt{remove()} видаляє з колекції \textbf{один конкретний елемент} (за його індексом або за значенням) і зменшує розмір структури на одиницю. Метод \texttt{clear()} повністю очищає колекцію, видаляючи абсолютно всі елементи одночасно й скидаючи розмір до нуля (\texttt{size = 0}).

    \item[10. Для чого використовується метод \texttt{contains()}?] ~\\
    Метод \texttt{contains(Object o)} використовується для перевірки наявності конкретного об'єкта всередині колекції. Він запускає внутрішній цикл порівняння елементів і повертає логічне значення \texttt{true}, якщо об'єкт знайдено, або \texttt{false}, якщо він відсутній.

    \item[11. Що повертає метод \texttt{size()}?] ~\\
    Метод \texttt{size()} повертає поточну кількість елементів (або пар «ключ -- значення» для \texttt{Map}), які фактично містяться в колекції на момент виклику. Повертає значення примітивного типу \texttt{int}.

    \item[12. Для чого використовується \texttt{isEmpty()}?] ~\\
    Метод \texttt{isEmpty()} призначений для швидкої перевірки колекції на наявність елементів. Він повертає логічне значення \texttt{true}, якщо в колекції немає жодного елемента (тобто \texttt{size() == 0}), або \texttt{false}, якщо колекція містить хоча б один об'єкт.

    \item[13. Що таке \texttt{Set}?] ~\\
    \texttt{Set} (множина) — це інтерфейс \textit{Java Collections Framework}, який описує колекцію, призначену для зберігання виключно \textbf{унікальних елементів}. На відміну від списків, множини не підтримують індексацію елементів, а порядок їхнього збереження (для \texttt{HashSet}) не гарантується.

    \item[14. Чим \texttt{Set} відрізняється від \texttt{List}?] ~\\
    Головна відмінність полягає в унікальності: \texttt{List} дозволяє зберігати дублікати й суворо упорядковує елементи за індексами, тоді як \texttt{Set} повністю блокує появу однакових значень. Крім того, у списку можна отримати елемент за номером за константний час через \texttt{get(i)}, а в множині пошук здійснюється ітераційно або за хеш-кодом.

    \item[15. Чому \texttt{HashSet} не допускає дублювання елементів?] ~\\
    Внутрішній устрій \texttt{HashSet} базується на роботі асоціативної карти \texttt{HashMap}, де додані елементи стають ключами. При спробі додати елемент, \texttt{HashSet} розраховує його \texttt{hashCode()} та перевіряє вміст через метод \texttt{equals()}. Якщо такий об'єкт уже існує, нове значення просто перезаписує старе за тим самим хеш-індексом, запобігаючи подвоєнню.

    \item[16. Що таке \texttt{Map}?] ~\\
    \texttt{Map} (карта/відображення) — це інтерфейс, який представляє структуру даних для зберігання елементів у вигляді **пар «ключ -- значення»**. \texttt{Map} не успадковує базовий інтерфейс \texttt{Collection}, оскільки оперує унікальними ключами, кожному з яких відповідає строго одне певне об'єктне значення.

    \item[17. Що таке ключ і значення у структурі \texttt{Map}?] ~\\
    \begin{itemize}
        \item \textbf{Ключ (\textit{Key}):} унікальний об'єктний ідентифікатор, який служить «первинним ключем» (як у базах даних) для миттєвого пошуку пов'язаного об'єкта. Дублікати ключів заборонені.
        \item \textbf{Значення (\textit{Value}):} дані (об'єкт), які прив'язані до конкретного ключа і витягуються за його запитом. Значення можуть дублюватися в різних ключах.
    \end{itemize}

    \item[18. Для чого використовується метод \texttt{put()}?] ~\\
    Метод \texttt{put(K key, V value)} призначений для додавання нової пари «ключ -- значення» до структури \texttt{Map}. Якщо операція виконується успішно, метод повертає \texttt{null}.

    \item[19. Що відбувається при додаванні нового значення за вже існуючим ключем?] ~\\
    Оскільки ключі в \texttt{Map} зобов'язані бути унікальними, при спробі викликати \texttt{put()} з уже існуючим у таблиці ключем, старе об'єктне значення за цим індексом буде **повністю стерте і замінене новим** значенням, що передане як аргумент. При цьому метод \texttt{put()} поверне посилання на те старе значення, яке щойно було видалене.

    \item[20. Для чого використовується \texttt{containsKey()}?] ~\\
    Метод \texttt{containsKey(Object key)} використовується для швидкої перевірки наявності конкретного ключа-індексу в асоціативній карті. Він повертає \texttt{true}, якщо такий ключ зареєстрований, або \texttt{false}, якщо ключ відсутній, що дозволяє уникнути помилки NullPointerException перед викликом \texttt{get()}.

    \item[21. Для чого використовується \texttt{entrySet()}?] ~\\
    Метод \texttt{entrySet()} використовується для отримання динамічної множини всіх записів відображення у вигляді \texttt{Set<Map.Entry<K, V>>}. Кожен елемент такої множини є повноцінним об'єктом пари, що дозволяє в одному циклі \texttt{for-each} одночасно й високоефективно витягувати як ключ через \texttt{getKey()}, так і пов'язане з ним значення через \texttt{getValue()}.

    \item[22. Що таке \texttt{generics}?] ~\\
    \texttt{Generics} (узагальнення / параметризовані типи) — це спеціальний архітектурний механізм мови Java, який дозволяє проектувати класи, інтерфейси та методи з параметрами, що виступають у ролі типів даних. Конкретний тип об'єктів підставляється розробником безпосередньо під час оголошення екземпляра.

    \item[23. Для чого використовуються \texttt{generics} під час роботи з колекціями?] ~\\
    Узагальнення використовуються для забезпечення \textbf{суворої типобезпеки (\textit{Type Safety})} на етапі компіляції додатка. Вони гарантують, що в колекцію не зможе потрапити сторонній об'єкт невідповідного типу, а також повністю позбавляють розробника необхідності виконувати небезпечне ручне приведення типів даних при витягуванні об'єктів із контейнера.

    \item[24. Що означає конструкція \texttt{List<String>}?] ~\\
    Дана конструкція означає параметризований динамічний список, який жорстко обмежений на роботу **виключно з об'єктами класу \texttt{String}**. Спроба додати туди число чи будь-який інший тип даних буде негайно заблокована компілятором Java, видавши помилку синтаксису.

    \item[25. Що означає конструкція \texttt{Map<Integer, Student>}?] ~\\
    Ця конструкція описує структуру асоціативної карти (відображення), в якій як унікальні ключі-ідентифікатори використовуються об'єкти класу-обгортки \texttt{Integer} (наприклад, ID або код студента), а як пов'язані з ними значення зберігаються реальні об'єкти користувацького класу \texttt{Student}.

    \item[26. У яких випадках доцільно використовувати \texttt{List}?] ~\\
    Списки \texttt{List} доцільно використовувати тоді, коли для програми критично важливо суворо **зберігати початковий порядок додавання елементів**, мати можливість швидкого доступу до об'єктів за їхніми порядковими номерами (індексами), а також коли бізнес-логіка повністю допускає наявність повторюваних значень (дублікатів).

    \item[27. У яких випадках доцільно використовувати \texttt{Set}?] ~\\
    Множини \texttt{Set} є незамінними в задачах, де ключовою вимогою є **гарантія унікальності елементів**. Їх доцільно обирати для фільтрації даних, виключення появи дублікатів (наприклад, формування унікального переліку жанрів кіно чи назв категорій товарів) та швидкої перевірки наявності об'єкта через метод \texttt{contains()}.

    \item[28. У яких випадках доцільно використовувати \texttt{Map}?] ~\\
    Структуру \texttt{Map} доцільно використовувати у випадках, коли необхідно організувати швидкий пошук об'єктів за принципом **«первинного ключа» (Primary Key)**. Вона ідеально підходить для проектування індексних таблиць, телефонних довідників, систем авторизації або логічного зв'язування унікальних ідентифікаторів (ID) з великими об'єктами сутностей.
\end{description}
